\documentclass[11pt, a4paper]{report}
\usepackage[utf8]{inputenc}
\usepackage[T1]{fontenc}
\usepackage[french]{babel}
\usepackage{amsmath, amssymb, amsfonts}
\usepackage{listings}
\usepackage{xcolor}
\usepackage{geometry}
\geometry{margin=2.5cm}

% --- Syntax Highlighting for Python ---
\definecolor{codegreen}{rgb}{0,0.6,0}
\definecolor{codegray}{rgb}{0.5,0.5,0.5}
\definecolor{codepurple}{rgb}{0.58,0,0.82}
\definecolor{backcolour}{rgb}{0.95,0.95,0.92}

\lstset{
    backgroundcolor=\color{backcolour},   
    commentstyle=\color{codegreen},
    keywordstyle=\color{magenta},
    numberstyle=\tiny\color{codegray},
    stringstyle=\color{codepurple},
    basicstyle=\ttfamily\small,
    breakatwhitespace=false,         
    breaklines=true,                 
    captionpos=b,                    
    keepspaces=true,                 
    numbers=left,                    
    numbersep=5pt,                  
    showspaces=false,                
    showstringspaces=false,
    showtabs=false,                  
    tabsize=2
}

\begin{document}

\section{Formalisme de l'Animation : Dynamique des Paquets d'Ondes}

Cette note détaille la structure mathématique du script \texttt{wavepacket\_animation.py}. Elle explicite le passage de l'équation du mouvement discrète à la simulation instationnaire par décomposition modale.

\subsection{Symétrisation de l'Opérateur Dynamique}

Pour une chaîne de $N$ masses, l'équation du mouvement s'écrit sous forme matricielle :
\begin{equation}
    \mathbf{M} \mathbf{\ddot{u}} + \mathbf{K} \mathbf{u} = 0
\end{equation}
où $\mathbf{M} = \text{diag}(m_n)$ et $\mathbf{K}$ est la matrice de raideur tridiagonale. Pour utiliser les solveurs de valeurs propres optimisés (\texttt{eigh}), nous devons transformer ce problème en un système symétrique standard.

\paragraph{Changement de variable :} 
On pose $\mathbf{v} = \mathbf{M}^{1/2} \mathbf{u}$. L'équation devient :
\begin{equation}
    \mathbf{\ddot{v}} + \mathbf{D} \mathbf{v} = 0, \quad \text{avec } \mathbf{D} = \mathbf{M}^{-1/2} \mathbf{K} \mathbf{M}^{-1/2}
\end{equation}
La matrice dynamique $\mathbf{D}$ est symétrique. Ses éléments sont :
\begin{equation}
    D_{n,n} = \frac{K_n + K_{n-1}}{m_n}, \quad D_{n,n+1} = -\frac{K_n}{\sqrt{m_n m_{n+1}}}
\end{equation}

\paragraph{Implémentation Python :}
\begin{lstlisting}[language=Python]
# Construction de la matrice D symetrique
for i in range(N):
    D[i, i] = (kl + kr) / m[i]
    if i < N-1:
        off = -K / np.sqrt(m[i] * m[i+1])
        D[i, i+1] = D[i+1, i] = off
\end{lstlisting}

\subsection{Évolution Temporelle par Décomposition Modale}

La solution générale du système est une superposition des modes propres $\mathbf{\Psi}_k$ oscillant aux pulsations $\omega_k$. Pour une condition initiale sans vitesse ($\mathbf{\dot{u}}(0) = 0$), le déplacement au temps $t$ est :
\begin{equation}
    \mathbf{v}(t) = \sum_{k=1}^N \alpha_k \mathbf{\Psi}_k \cos(\omega_k t)
\end{equation}

\paragraph{Projection Spectrale :}
Les coefficients $\alpha_k$ sont obtenus en projetant la condition initiale $\mathbf{v}(0)$ sur la base des vecteurs propres :
\begin{equation}
    \boldsymbol{\alpha} = \mathbf{\Psi}^T \mathbf{v}(0)
\end{equation}

\paragraph{Implémentation Python :}
\begin{lstlisting}[language=Python]
# Resolution spectrale
omega2, Psi = eigh(D)
omega = np.sqrt(np.clip(omega2, 1e-9, None))

# Projection et evolution
coeffs = Psi.T @ v0
v_t = Psi @ (coeffs * np.cos(omega * t))
\end{lstlisting}

\subsection{Synthèse du Paquet d'Ondes}

La condition initiale est un paquet d'ondes gaussien localisé autour d'un atome central $n_c$. Sa forme mathématique est définie par le produit d'une enveloppe et d'une porteuse :
\begin{equation}
    u_n(0) = \exp\left( -\frac{(n-n_c)^2}{2\sigma^2} \right) \cdot \cos(q_0 (n-n_c))
\end{equation}
où $\sigma$ dicte l'étalement spatial et $q_0$ est le vecteur d'onde central (fréquence spatiale).

\paragraph{Implémentation Python :}
\begin{lstlisting}[language=Python]
# Definition de l'etat initial
u0 = np.exp(-(X - center)**2 / (2 * sigma**2))
if q0 > 0:
    u0 *= np.cos(q0 * (X - center))
u0 /= np.linalg.norm(u0) # Normalisation de l'energie totale
\end{lstlisting}

\end{document}
